\chapter{模型与模拟方法}

\section{格点气体模型}

考虑一个$N\times N$的正方格点，采用开放边界条件。每个格点$(i,j)$（$i,j=1,2,\ldots,N$）上定义一个
占据变量$n_{ij}\in\{0,1\}$：$n_{ij}=1$表示该格点放置了一个皇后，$n_{ij}=0$表示该格点为空。
系统的哈密顿量（能量函数）计数所有互相攻击的皇后对数：
%
\begin{equation}
  H = J\sum_{\langle(i,j),(i',j')\rangle_{\text{attack}}} n_{ij}\,n_{i'j'},
  \label{eq:H}
\end{equation}
%
其中求和遍历所有共享同一行（$i=i'$）、同一列（$j=j'$）、同一主对角线（$i-j=i'-j'$）
或同一反对角线（$i+j=i'+j'$）的皇后对。耦合常数$J>0$确保排斥相互作用，全文设定$J=1$作为能量单位。
基态（$H=0$）精确对应于所有皇后互不攻击的构型——即N皇后问题的解。因此，基态简并度就是$Q(N)$。

哈密顿量的攻击线分解提供了更简洁的形式。将棋盘分解为攻击线的集合$\mathcal{L}$：
$N$条行、$N$条列、$2N-1$条主对角线、$2N-1$条反对角线，共计$6N-2$条攻击线。
设$n_{\alpha}$为第$\alpha$条攻击线上的皇后数，则该攻击线内部的攻击对数为组合数
$\binom{n_{\alpha}}{2}=n_{\alpha}(n_{\alpha}-1)/2$。由此，哈密顿量可重写为所有攻击线上
攻击对数的总和：
%
\begin{equation}
  H = \sum_{\alpha\in\mathcal{L}} \binom{n_{\alpha}}{2}
    = \sum_{\alpha\in\mathcal{L}} \frac{n_{\alpha}(n_{\alpha}-1)}{2}.
  \label{eq:Hlines}
\end{equation}
%
这一分解形式将总能量表达为$6N-2$个独立攻击线的贡献之和，大大简化了高温极限下
能量期望值的计算。

\begin{figure}[htbp]
  \centering
  \includegraphics[width=0.55\textwidth]{fig1_schematic.png}
  \caption{$N=8$棋盘上的一个非攻击基态构型示例（$Q(8)=92$个解之一）。无任何两个皇后
  共享同一行、列或对角线。}
  \label{fig:schematic}
\end{figure}

物理上，攻击线分解直观地反映了约束的几何结构。行和列分别对应水平和垂直方向上的
$N$条等长直线，每条线的长度均为$N$。主对角线方向（$i-j=\mathrm{const}$）
和反对角线方向（$i+j=\mathrm{const}$）则各包含$2N-1$条不等长的线：
长度从1（棋盘角落的单个格点）递增到$N$（中心贯穿棋盘的主对角线），再递减回1。
攻击线长度的不均匀分布是两族对角线约束效率低于行/列约束的根本原因。

对于N皇后问题，一个隐含的约束是每行恰好放置一个皇后——这是由皇后总数为$N$、
行为$N$条这一天然对应关系所决定的。在我们的模拟中，我们选择在正则系综中工作：
固定皇后总数为$N$，模拟温度$T$下的热平衡态。

\section{蒙特卡洛模拟方法}

\subsection{Kawasaki动力学与Metropolis抽样}

我们采用Kawasaki动力学（粒子-空位交换）在固定皇后数$N$的正则系综中进行
Monte Carlo抽样。每一步尝试的具体操作如下：从当前构型中随机选择一个被皇后占据的
格点和任意一个空位格点，提议交换这两个格点的占据状态（即皇后移动到空位上）。
提议的接受概率由Metropolis准则决定：
%
\begin{equation}
  P_{\text{accept}} = \min\left(1, \exp\left(-\Delta E / T\right)\right),
  \label{eq:metropolis}
\end{equation}
%
其中$\Delta E$是提议移动引起的能量变化，$T$是温度（$k_{B}=1$，全文采用无量纲化方案）。
一次``sweep''定义为$N$次交换尝试——平均每个皇后被提议移动一次。

选择Kawasaki动力学（皇后-空位交换）而非皇后-皇后交换的原因在于：
前者允许单个皇后通过一系列交换移动穿越整个棋盘，具有遍历整个配置空间的充分能力；
而皇后-皇后交换仅置换两个已有皇后的位置，无法改变哪些格点被占据这一全局信息，
因而不能有效探索具有不同空间结构的构型区域。

\subsection{模拟参数与温度网格}

我们对$N=8,16,32,64,128,256,512,1024$共8个系统尺寸分别进行了模拟。
每个温度点的测量阶段进行$10^{8}$次sweep，热化阶段进行$2\times10^{6}$次sweep
（远大于最大自相关时间，确保充分平衡）。温度网格在$T=0.05\,J$到$500\,J$之间
均匀布置了280个温度点：在比热峰区域（$T=0.1025\sim0.500\,J$，步长$0.0025\,J$）
布置了160个密集点以保证数值积分精度；低温区（$T=0.050\sim0.100\,J$，步长
$0.005\,J$）布置11个点；其余109个点分布在中、高温区。所有8个系统尺寸使用
完全相同的温度网格和测量参数，保证了数据集的自洽性。

\subsection{物理量测量与误差分析}

在每个温度点、热化完成后，我们在$10^{8}$次测量sweep中记录每一步的能量$E$。
从这些时序数据中计算：
%
\begin{align}
  \langle E\rangle &= \frac{1}{M}\sum_{t=1}^{M} E(t), \\
  \langle E^{2}\rangle &= \frac{1}{M}\sum_{t=1}^{M} E^{2}(t), \\
  C_{v}/N &= \frac{\langle E^{2}\rangle - \langle E\rangle^{2}}{T^{2}N}.
\end{align}

统计误差的估计采用jackknife重采样方法\cite{Efron1982}。将全部$10^{8}$次测量
sweep分成200个等长的bin（每个bin包含$5\times10^{5}$次sweep），
通过删除第$k$个bin后重新计算物理量，利用200个折刀估计值之间的涨落来估计标准误差。
我们对不同的随机初始构型进行了独立模拟，所有独立运行在全部温度点上产生了统计一致的
结果，确认了Kawasaki动力学在给定热化和测量窗口内能够充分遍历构型空间。

\subsection{并行计算策略与收敛性说明}

模拟在280个CPU核心上并行执行，每个核心负责一个温度点的所有8个系统尺寸的串行模拟。
整个计算任务的总墙钟时间约为9小时。

值得一提的是，N皇后格点气体的模拟在计算上具有一定特殊挑战。由于低温下接受率极低
（$T<0.1\,J$时接受率$<10^{-5}$），系统在基态构型附近近乎冻结，
自相关时间在$T\approx0.055\sim0.10\,J$处出现尖锐峰值
（$N=1024$时$\tau_{\text{int}}\approx5.0\times10^{4}$ sweeps）。
然而，关键的物理特征是比热峰位于$T^{*}\approx0.235\,J$——远高于自相关时间的
峰值温度。在$T^{*}$处，自相关时间仅为$85\sim200$ sweeps，提供超过
$5\times10^{5}$个独立样本，确保了比热峰区域的极高统计精度。
低温区虽自相关性较强，但能量和比热本身在此温度已呈指数衰减趋于零，
对热力学积分结果的贡献微乎其微。

\subsection{模型验证}

为验证模拟程序的正确性，我们在$N=8$和$N=16$两个尺寸上进行了详细测试。
对于$N=8$，$Q(8)=92$是精确已知的；对于$N=16$，$Q(16)=14{,}772{,}512$
同样是已知的精确值。我们在极低温度验证了热化后的能量趋向于零、
系统可以在不同的基态构型之间随机游走，以及通过热力学积分反推出的基态熵
$s_{0}^{\text{MC}}$与精确值$\ln Q(N)/N$的偏差小于$0.1\%$，
验证了整个模拟$+$积分流程的正确性。
